\subsubsection{Symbian OS}

\sourcenotebox{\textbf{Source Note:} The listings below are taken from the
local Symbian EPOC Kernel Architecture 2 (EKA2) source tree in this workspace
under \texttt{kernel/eka/}.
The upstream repository is
\url{https://github.com/SymbianSource/oss.FCL.sf.os.kernelhwsrv}. Each caption
uses the shortest unambiguous path under \texttt{kernel/eka/}.}

Symbian EKA2 treats memory management as several cooperating subsystems rather
than one allocator. The code in \texttt{kernel/eka/} separates physical-page
metadata, per-process virtual address spaces, process-visible mappings, demand
paging, random-access memory (RAM) zoning, and recovery from fragmentation or
pressure \citep{sales2005, symbian_eka2_kernel}. That division is what makes
the memory behavior source-legible: most major policies are stated directly in
types, creation paths, and paging interfaces.

\paragraph{Memory Model and Address Spaces}

The common memory-management-unit (MMU) base code already reveals the core
model. Listing~\ref{lst:mem-pageinfo-addressspace} shows that each physical
page carries explicit ownership and paging-state metadata, while each process
gets a dedicated address-space object with its own page directory, virtual
allocator, and mapping list. The comment also makes the operating-system
address space identifier (OS ASID) policy explicit: the kernel refers to
address spaces by indexed IDs rather than by raw pointers in most interfaces.

\begin{listing}[H]
\begin{minted}{cpp}
struct SPageInfo
	{
	enum TType
		{
		EInvalid=0,
		EUnused=2,
		EChunk=3,
		EPageDir=7,
		EPagedCode=11,
		EPagedFree=14,
		};

	enum TState
		{
		EStateNormal = 0,
		EStatePagedYoung = 1,
		EStatePagedOld = 2,
		EStatePagedDead = 3,
		EStatePagedLocked = 4
		};
	};

class DAddressSpace : public DReferenceCountedObject
	{
public:
	static TInt New(TPhysAddr& aPageDirectory);
	};
\end{minted}
\caption{Physical-page metadata and per-process address-space objects. Files:
\texttt{mmubase.h} and \texttt{flexible/mmu/maddressspace.h}.}
\label{lst:mem-pageinfo-addressspace}
\end{listing}

\paragraph{Process Setup and Paging Policy}

Address-space creation happens in the normal process constructor rather than in
a later memory-specific pass. Listing~\ref{lst:mem-process-create} shows that
kernel processes stay in the kernel address space, while user processes call
\texttt{MM::AddressSpaceAlloc()} and store the returned OS ASID. The same setup
path also decides whether process data should be paged, so paging policy is
bound to process creation and image flags instead of being bolted on after the
fact.

\begin{listing}[H]
\begin{minted}{cpp}
TInt DMemModelProcess::SetPaging(const TProcessCreateInfo& aInfo)
	{
	TUint pagedFlags = aInfo.iFlags & TProcessCreateInfo::EDataPagingMask;
	if (pagedFlags == TProcessCreateInfo::EDataPagingMask)
		{
		return KErrCorrupt;
		}
	if (aInfo.iAttr & ECodeSegAttKernel)
		{
		return KErrNone;
		}
	/* otherwise follow the configured data-paging policy */
	TUint dataPolicy = TheSuperPage().KernelConfigFlags()
		& EKernelConfigDataPagingPolicyMask;
	if (dataPolicy == EKernelConfigDataPagingPolicyAlwaysPage)
		{
		iAttributes |= EDataPaged;
		return KErrNone;
		}
	if (pagedFlags == TProcessCreateInfo::EDataPaged)
		{
		iAttributes |= EDataPaged;
		}
	return KErrNone;
	}

TInt DMemModelProcess::DoCreate(TBool aKernelProcess, TProcessCreateInfo& aInfo)
	{
	TInt r=KErrNone;

	if (aKernelProcess)
		{
		iAttributes |= ESupervisor;
		iOsAsid = KKernelOsAsid;
		}
	else
		{
		r = MM::AddressSpaceAlloc(iPageDir);
		if (r>=0)
			{
			iOsAsid = r;
			r = KErrNone;
			}
		}
	if (r == KErrNone)
		{
		__e32_atomic_store_ord32(&iOsAsidRefCount, 1);
		}
	return r;
	}
\end{minted}
\caption{Process address-space allocation and data-paging policy. File:
\texttt{flexible/mprocess.cpp}.}
\label{lst:mem-process-create}
\end{listing}

\paragraph{Data, BSS, Stack, and Code Regions}

Process memory is assembled from explicit kernel memory objects and mappings,
not from a single anonymous heap. Listing~\ref{lst:mem-data-stack} shows two
variants of the same idea. In the flexible model, Symbian rounds the combined
data and block-started-by-symbol (BSS) size to pages, allocates backing memory,
and creates a user mapping for it. In the multiple model, it instead creates a
disconnected user-data chunk sized for static data plus future stack growth.

\begin{listing}[H]
\begin{minted}{cpp}
TInt DMemModelProcess::CreateDataBssStackArea(TProcessCreateInfo& aInfo)
	{
	TInt r = KErrNone;
	TInt dataBssSize = MM::RoundToPageSize(aInfo.iTotalDataSize);
	if(dataBssSize)
		{
		DMemoryObject* memory;
		TMemoryObjectType memoryType = iAttributes&EDataPaged
			? EMemoryObjectPaged
			: EMemoryObjectMovable;
		r = MM::MemoryNew(memory,memoryType,MM::BytesToPages(dataBssSize));
		if(r==KErrNone)
			{
			/* allocate pages, map them into the process, record the base */
			r = MM::MemoryAlloc(memory,0,MM::BytesToPages(dataBssSize));
			r = MM::MappingNew(iDataBssMapping,memory,EUserReadWrite,OsAsid());
			iDataBssRunAddress = MM::MappingBase(iDataBssMapping);
			}
		}
	return r;
	}

TInt DMemModelProcess::CreateDataBssStackArea(TProcessCreateInfo& aInfo)
	{
	TInt dataBssSize=Mmu::RoundToPageSize(aInfo.iTotalDataSize);
	TInt maxSize=dataBssSize+PP::MaxStackSpacePerProcess;

	SChunkCreateInfo cinfo;
	cinfo.iGlobal=EFalse;
	cinfo.iAtt=TChunkCreate::EDisconnected;
	cinfo.iOperations=SChunkCreateInfo::EAdjust|SChunkCreateInfo::EAdd;
	cinfo.iType=EUserData;
	cinfo.iMaxSize=maxSize;
	cinfo.iInitialTop=dataBssSize;
	cinfo.iName.Set(KDollarDat);
	cinfo.iOwner=this;
	TLinAddr cb;
	return NewChunk((DChunk*&)iDataBssStackChunk,cinfo,cb);
	}
\end{minted}
\caption{Static-data and stack-region construction in the flexible and
multiple memory models. Files: \texttt{flexible/mprocess.cpp} and
\texttt{multiple/mprocess.cpp}.}
\label{lst:mem-data-stack}
\end{listing}

Executable attachment then fixes the actual run addresses for code and data.
Listing~\ref{lst:mem-codeseg-attach} shows that the process pulls those
addresses from the code segment metadata, using one path for execute-in-place
images and another for RAM-loaded images. The important point is that code,
data, and export metadata all become part of process creation, not a later
loader patch-up step outside the kernel memory model.

\begin{listing}[H]
\begin{minted}{cpp}
TInt DEpocProcess::AttachExistingCodeSeg(TProcessCreateInfo& aInfo)
	{
	DEpocCodeSeg& s=*(DEpocCodeSeg*)iTempCodeSeg;
	if (s.iAttachProcess && s.iAttachProcess!=this)
		return KErrAlreadyExists;
	TInt r=CreateDataBssStackArea(aInfo);
	if (r==KErrNone)
		{
		if (s.iXIP)
			{
			const TRomImageHeader& rih=s.RomInfo();
			aInfo.iCodeLoadAddress=rih.iCodeAddress;
			aInfo.iCodeRunAddress=rih.iCodeAddress;
			aInfo.iDataRunAddress=rih.iDataBssLinearBase;
			}
		else
			{
			SRamCodeInfo& ri=s.RamInfo();
			aInfo.iCodeRunAddress=ri.iCodeRunAddr;
			aInfo.iDataRunAddress=ri.iDataRunAddr;
			}
		}
	return r;
	}
\end{minted}
\caption{Executable attachment and run-address assignment. File:
\texttt{pprocess.cpp}.}
\label{lst:mem-codeseg-attach}
\end{listing}

\paragraph{Demand Paging and Reclaim}

Demand paging is not implicit background magic; the intended policy is
documented in the headers. Listing~\ref{lst:mem-demand-paging} describes a live
page list split into young and old pages, where faults on old pages rejuvenate
them into the young list. The same listing also shows the pager interface for
handling a fault, freeing cache pages for the rest of the system, flushing
paged regions, and reclaiming previously donated pages. This is a pseudo Most
Recently Used (MRU) design rather than an exact least recently used policy.

\begin{listing}[H]
\begin{minted}{cpp}
/**
Young pages are accessible; old pages fault and are moved back to the young
list. If no free RAM page is available, the oldest live page is recycled.
*/
class DemandPaging : public RamCacheBase
	{
public:
	virtual TInt Fault(TAny* aExceptionInfo)=0;
	};

class DPager
	{
public:
	TBool GetFreePages(TInt aNumPages);
	TInt HandlePageFault(/* fault context omitted */);
	void FlushAll();
	TInt FlushRegion(DMemModelProcess* aProcess, TLinAddr aStartAddress, TUint aSize);
	void DonatePages(TUint aCount, TPhysAddr* aPages);
	TInt ReclaimPages(TUint aCount, TPhysAddr* aPages);
	};
\end{minted}
\caption{Demand-paging policy and pager control surface. Files:
\texttt{demand\_paging.h} and \texttt{flexible/mmu/mpager.h}.}
\label{lst:mem-demand-paging}
\end{listing}

\paragraph{Physical RAM and Kernel Allocation}

Below the pager, physical memory is zone-aware. Listing~\ref{lst:mem-ram-heap}
shows the zone-counting path validating alignment, uniqueness, and full
coverage of allocatable RAM before the zones can be used. Kernel
dynamic allocation is handled separately by \texttt{RHeapK}, which serializes
access with a mutex and leaves compression disabled. That split keeps page
ownership and heap allocation distinct instead of conflating them.

\begin{listing}[H]
\begin{minted}{cpp}
void DRamAllocator::CountZones(const SRamZone* aZones)
	{
	TUint32 totalSize = 0;
	const SRamZone* pCurZ = aZones;
	for (; pCurZ->iSize != 0; pCurZ++)
		{
		TUint32 curEnd = pCurZ->iBase + pCurZ->iSize - 1;
		/* alignment and overlap checks omitted */
		if (curEnd <= pCurZ->iBase)
			{
			Panic(EZonesAlignment);
			}
		if (pCurZ->iId == KRamZoneInvalidId)
			{
			Panic(EZonesIDInvalid);
			}
		iNumZones++;
		totalSize += pCurZ->iSize;
		}
	if (totalSize>>KPageShift < iTotalRamPages)
		{
		Panic(EZonesIncomplete);
		}
	}

RHeapK::RHeapK(TInt aInitialSize)
	: RHybridHeap(aInitialSize, 0, EFalse)
	{
	}

TInt RHeapK::Compress()
	{
	return 0;
	}

TInt RHeapK::CreateMutex()
	{
	DMutex*& m = *(DMutex**)&iLock;
	return K::MutexCreate(m, KLitKernHeap, NULL, EFalse, KMutexOrdKernelHeap);
	}
\end{minted}
\caption{Zone-aware RAM setup and the kernel heap path. Files:
\texttt{ramalloc.cpp} and \texttt{kheap.cpp}.}
\label{lst:mem-ram-heap}
\end{listing}

\paragraph{Protection and Paging Exceptions}

Paging faults are isolated from ordinary execution with a dedicated trap type.
Listing~\ref{lst:mem-paging-trap} shows that pageable operations can be wrapped
in \texttt{XTRAP\_PAGING}, which installs a \texttt{TPagingExcTrap} before the
memory access and removes it afterward. That gives the kernel a separate
control path for recoverable paging exceptions instead of treating every fault
as a fatal access violation.

\begin{listing}[H]
\begin{minted}{cpp}
class TPagingExcTrap
	{
public:
	TInt Trap();
	inline void UnTrap();
	void Exception(TInt aResult);
public:
	TUint32 iState[EXC_TRAP_CTX_SZ];
	DThread* iThread;
	};

#define XTRAP_PAGING(_r,_s)	{								\
							TPagingExcTrap __t;				\
							if (((_r)=__t.Trap())==0)		\
								{							\
								_s;							\
								__t.UnTrap();				\
								}							\
							}
\end{minted}
\caption{Paging-exception trapping for pageable kernel accesses. File:
\texttt{kern\_priv.h}.}
\label{lst:mem-paging-trap}
\end{listing}

\paragraph{Defragmentation and Memory Pressure}

Symbian does not leave fragmentation and pressure to passive allocator luck.
Listing~\ref{lst:mem-defrag-thrash} shows two active control paths. The
defragmenter is initialized on its own deferred function call (DFC) queue and
can walk a zone, moving movable pages elsewhere to recover better physical
contiguity. Separately, \texttt{DThrashMonitor} measures whether paging itself
is becoming pathological. Together with the pager interface in
Listing~\ref{lst:mem-demand-paging}, this shows that memory pressure is managed
as an ongoing kernel policy problem, not just a one-shot allocation failure.

\begin{listing}[H]
\begin{minted}{cpp}
void Defrag::Init3(DRamAllocator* aRamAllocator)
	{
	TheDefrag = this;

	TInt r = Kern::DfcQInit(&iTaskQ, KDefragIdlingThreadPriority,
		&KDefragDfcThreadName);
	if (r!=KErrNone)
		Panic(EDfcQInitFailed);

	iRamAllocator = aRamAllocator;
	}

TInt Defrag::ClearMovableFromZone(SZone& aZone, TBool aBestEffort,
	TRamDefragRequest* aRequest)
	{
	TPhysAddr newAddr;
	const TUint moveDisFlags = (aBestEffort)
		? 0
		: M::EMoveDisBlockRest | M::EMoveDisMoveDirty;
	TUint offset = 0;
	TInt ret = iRamAllocator->NextAllocatedPage(&aZone, offset, EPageMovable);
	while (ret == KErrNone)
		{
		if (aRequest->PollForCancel())
			{
			return KErrCancel;
			}
		/* move one movable page out of the zone */
		TPhysAddr addr = (offset << M::PageShift()) + aZone.iPhysBase;
		TInt moved = M::MovePage(addr, newAddr, aZone.iId, moveDisFlags);
		if (moved != KErrNone && !aBestEffort)
			{
			return moved;
			}
		offset++;
		ret = iRamAllocator->NextAllocatedPage(&aZone, offset, EPageMovable);
		}
	return KErrNone;
	}

class DThrashMonitor
	{
public:
	DThrashMonitor();
	void Start();
	TInt ThrashLevel();
	TInt SetThresholds(TUint thrashing, TUint good);
	void NotifyStartPaging();
	void NotifyEndPaging();
	};
\end{minted}
\caption{Active defragmentation and thrash monitoring. Files:
\texttt{defragbase.cpp} and \texttt{flexible/mmu/mthrash.h}.}
\label{lst:mem-defrag-thrash}
\end{listing}
